% SPDX-License-Identifier: CC-BY-4.0
\section{Conclusion}
\label{sec:conclusion}

Behavioral host intrusion detection systems all rest on an
assumption about their learning window that real-world adversary
dwell times invalidate. We formalized the resulting baseline-poisoning
threat model, demonstrated that the naive estimators in widespread
use collapse under modest poisoning budgets, proposed \sediment{} as
a per-feature trimmed-moment estimator with a closed-form
Gaussian-truncation bias correction, and evaluated it across nine
experiments including a white-box adversary that knows the
defender's hyperparameters. \sediment{} provides a defensible
robustness guarantee at a transparent operational cost: a small but
quantifiable bump in clean-condition false-positive rate when the
benign distribution leaves the Gaussian null, and effectively zero
cost when it does not. The white-box experiment gives evidence the
obvious attack class does not improve on the budget-bounded
black-box adversary the construction was designed against.

The artifact is open source. The evaluation regenerates from a
single command. The Python reference and Rust production
implementations are bit-equivalent through a CI-enforced
cross-implementation parity check. We hope the threat model and
evaluation harness are useful to others building behavioral
detectors at any layer of the stack.
